% Complete dense-order extension and corrected Kamp IV.2 construction.

\subsection{Dense-order separation and the integer unary basis}
\label{sec:annex-p09-kamp-iv2-integer-basis}

\begin{annexitem}{annex:p9.kamp.dense.local.unary}{thm:p09-kamp-dense-local-unary}{Finite partial maps for a convex interval}
\begin{proof}
Let \(u,v\in I\) have the same membership in \(A\).
We construct a family of finite, one-to-one partial maps from
\(T\) to itself containing the map \(u\mapsto v\).
Every map will preserve order and membership in \(A\), and
every finite map in the family will admit an extension including
any specified point in its domain, or in its range. We use these
extension properties directly to prove preservation of all
first-order formulas; an automorphism of the full order is not
required.

First note that any point outside \(I\) lies below all of
\(I\) or above all of \(I\). Otherwise it would lie
between two points of the convex set and would itself belong
to \(I\). We require every selected point outside \(I\)
to map to itself. This preserves its predicate value, all order
relations among such points, and all its order relations with
points of \(I\), whatever \(A\) does outside \(I\).

Suppose first that \(A\cap I\) is \(\varnothing\) or
\(I\). Within \(I\) use any finite strictly increasing
partial map containing \(u\mapsto v\). To include a new
domain point, look at its position relative to the finitely many
already selected points in \(I\), and choose an unused range
point in the corresponding position. Density supplies a point
between two bounds; absence of endpoints supplies a point when
only one bound exists or neither exists. A finite set of already
used points can be avoided by choosing in an open interval
between the nearest bounds. The same argument with domain and
range exchanged includes any prescribed new range point.
All predicate values inside \(I\) are the same, so these
extensions preserve the predicate. Together with identity outside
\(I\), they have all the required properties.

Suppose now that \(A\cap I\) is \(D\) or
\(I\setminus D\), with the two-sided markers \(d_n\).
Every point of \(I\) is a marker or lies in a unique interval
\((d_n,d_{n+1})\). Indeed, choose marker bounds
\(d_a<x<d_b\); only the finitely many indices from \(a\)
to \(b\) can determine its position between these bounds.
Selecting the last marker at or below \(x\) gives the claimed
alternative and uniqueness. Each open interval
\((d_n,d_{n+1})\) is nonempty, dense, and has no endpoints
in its induced order.

Since \(u,v\) have the same predicate value, either both
are markers or both are nonmarkers. If \(u=d_a,v=d_b\),
fix \(h=b-a\). If \(u\in(d_a,d_{a+1})\) and
\(v\in(d_b,d_{b+1})\), again fix \(h=b-a\).
Require selected markers to map by
\(d_n\mapsto d_{n+h}\), and selected points of
\((d_n,d_{n+1})\) to map increasingly into
\((d_{n+h},d_{n+h+1})\). These restrictions contain the
initial pair and preserve every order relation between distinct
marker intervals. They preserve the predicate: all markers have
one predicate value and all nonmarkers inside \(I\) the other.

For a new selected marker, use the prescribed translated marker.
No collision is possible, since different marker indices have
different translates and nonmarkers map to nonmarkers. For a
new point inside a marker interval, extend the finite increasing
map within that interval by density and absence of endpoints,
as in the constant-predicate case. For a point outside \(I\),
use identity. These choices preserve the prescribed restrictions
and all order and predicate relations. To include a new range
point, use the same constructions with the index translation
\(-h\) and the inverse partial map. Thus both extension
properties hold in this case as well.

We now prove the formula-preservation assertion. For any map in
the constructed family, and any finite tuple from its domain,
the tuple and its image agree on all atomic formulas: order is
preserved and so is the predicate. Boolean combinations preserve
agreement. For an existential formula, a witness on the domain
side can be included by the domain-extension property. The
induction hypothesis for its matrix gives a witness on the
range side. A witness on the range side is handled by the
range-extension property. Universal formulas follow by the
same argument in both directions, or by negated existential
quantification. This is structural induction over the complete
first-order formula, valid with quantifiers ranging over all
of \(T\). Applied to the initial pair, it says that \(u\)
and \(v\) satisfy the same one-free-variable formulas in
the stated signature.

In particular a defining formula for \(O\) has constant
truth value on \(A\cap I\) and on
\(I\setminus(A\cap I)\). Independently choosing whether
to include each of these two classes yields exactly the four
sets in the conclusion, with repetitions when one class is
empty. This proves the local closure assertion. The proof
uses density and the displayed marker decomposition, and makes
no completeness or cardinality assumption.
\end{proof}
\end{annexitem}

\begin{annexitem}{annex:p9.kamp.dense.unary.separation}{thm:p09-kamp-dense-unary-separation}{The separation inside every nontrivial dense order}
\begin{proof}
Choose \(a<b\) in \(T\). By density choose a point
\(d_0\) strictly between them. Recursively choose
\(d_{n+1}\) strictly between \(d_n\) and \(b\) for
\(n\geq0\), and \(d_{-n-1}\) strictly between \(a\)
and \(d_{-n}\) for \(n\geq0\). This gives an increasing
sequence indexed by all integers. Set
\[
\begin{aligned}
 I&=\{x\in T:\exists m,n\in\mathbb Z\ (d_m<x<d_n)\},\\
 E&=\{d_{2k}:k\in\mathbb Z\},\qquad
 H=\{d_{2k+1}:k\in\mathbb Z\}.
\end{aligned}
 \tag{1}
\]
The set \(I\) is nonempty and convex, and inherits density.
Every one of its points has a marker on each side, and markers
in turn have earlier and later markers, so \(I\) has no
endpoints. The sets \(E,H,E\cup H\) are each two-sided
marker sets inside \(I\). For \(E\) or \(H\), choose
even or odd indices beyond any given pair of marker bounds.
No limit of the chosen sequence is needed in \(T\).

Let \(\mathcal B_I\) be the field of subsets of \(I\)
generated by \(E,H\). Pairwise disjointness gives exactly
the possible forms
\[
 \mathcal B_I=
 \{\varnothing,I,E,H,E\cup H,I\setminus E,
                  I\setminus H,I\setminus(E\cup H)\}.
 \tag{2}
\]
For example, intersections of two unions of generators retain
the common indices; intersection with a complemented union
removes its indices; and the intersection of two complemented
unions is the complement of their union. These four cases show
that the displayed family is closed under complement and
intersection and contains the generators. Conversely every
displayed member is a Boolean combination of them. Thus (2)
holds. The third atom \(I\setminus(E\cup H)\) is nonempty
by density between consecutive markers.

Interpret \(q_0\) by \(E\) and \(q_1\) by \(H\)
as subsets of the whole of \(T\). We claim that for every
finite formula \(\eta\) in these two letters from
\(\mathsf{TM}_{\mathcal J}\),
\[
                       \llbracket\eta\rrbracket\cap I
                                    \in\mathcal B_I.
 \tag{3}
\]
Induct on its syntax. The atomic values are generators.
Intersection preserves the field, and global complement
restricts to complement inside \(I\). For the unary step,
put \(A=\llbracket\alpha\rrbracket\), where the induction
hypothesis gives \(A\cap I\in\mathcal B_I\). Every
possibility in (2) is constant on \(I\), a two-sided marker
set, or the complement of such a marker set. Hence
\cref{thm:p09-kamp-dense-local-unary} applies and puts
\(O(A)\cap I\) among
\(\varnothing,I,A\cap I,I\setminus(A\cap I)\), all
members of \(\mathcal B_I\). The operand \(A\) outside
\(I\) can have any value: the local lemma expressly permits
this, which is necessary after earlier unary operations.
Only this operand predicate is named in the defining formula
for \(O\); the other original input is not an additional
predicate parameter of that unary operation.

The same induction covers all finite truth-functional
connectives. Their output is a finite union of intersections
of input sets or their complements, with \(\varnothing\)
and \(I\) for the empty union and empty intersection.
This completes (3) for the full permitted alphabet.

Choose \(u\in(d_0,d_1)\) and
\(v\in(d_{-1},d_0)\). Both lie in the third atom of
\(\mathcal B_I\). Therefore every formula in (3) agrees
at \(u,v\). But \(S(E,T\setminus H)\) is true at
\(u\), witnessed by \(d_0\), and false at \(v\):
any earlier even marker is at most \(d_{-2}\), and its
interval to \(v\) contains \(d_{-1}\in H\).
Likewise \(U(E,T\setminus H)\) is true at \(v\),
witnessed by \(d_0\), and false at \(u\), because
every later even marker is at least \(d_2\), with
\(d_1\in H\) in its intervening interval.

If a two-input formula \(\eta\) defined either binary
operation, replace every \(q_1\) by \(\neg q_1\).
The resulting formula remains in the unary language. Its
value at the fixed inputs \((E,H)\) is the value of
\(\eta\) at \((E,T\setminus H)\): atomic substitution
has these two values, Boolean operations preserve them, and
each unary operation applies the same set function to equal
operand sets. This structural substitution argument proves
the equality for the whole finite formula. The resulting
truth set would distinguish \(u,v\), contrary to (3).
Each binary operation therefore yields a contradiction.

The construction stays inside \((a,b)\), so it also works
when \(T\) has a first point, a last point, or both. All
operator definitions were allowed to be specific to the
particular order. On a one-point order no strict witness
exists in either direction, and
\((q_0\land\neg q_0)\land q_1\)
expresses both constant-false operations while containing both
input letters. This verifies the
excluded boundary and establishes the source observation
with its exact dense-order hypotheses
\cite[ch.~IV, pp.~99--100]{Kamp1968}.
\end{proof}
\end{annexitem}

\begin{annexitem}{annex:p9.kamp.iv2.unary.semantics}{thm:p09-kamp-iv2-unary-semantics}{Deriving all four integer clauses}
\begin{proof}
An interval \((r,s)\cap\mathbb Z\), with \(r<s\),
is empty exactly when \(s=r+1\). If \(s\geq r+2\),
the integer \(r+1\) lies inside; if \(s=r+1\), no
integer lies strictly between them. A universal condition
requiring the contradiction \(\bot_{P_0}\) at every
interior point therefore forces this adjacency.

In \(\pi_-(i)\), the witness \(t_1<i\) must consequently
be \(i-1\), and its predicate requirement is exactly the
truth of the operand at \(i-1\). Conversely that truth
provides the witness and the universal clause is vacuous.
In \(\pi_+(i)\), adjacency instead forces \(t_1=i+1\),
and the identical argument yields the future-neighbor clause.
These are the claims for \(P_1,F_1\).

For \(\iota_-(i)\), choose its outer witness \(j<i\).
The inner contradictory universal forces \(t_2=j-1\),
and the predicate at \(t_2\) is required to be false.
The remaining universal says that the operand is false
throughout \((j,i)\cap\mathbb Z\). Together with its
truth at \(j\), these are precisely the three conditions
in the displayed clause for \(S_1\). Conversely, given
those conditions, choose \(t_1=j,t_2=j-1\); the inner
interval is empty and the outer interval condition is the
given one. For \(\iota_+(i)\), the same explicit
calculation forces \(t_1=j>i,t_2=j+1\), leaves the
operand false on \((i,j)\cap\mathbb Z\), and requires
its truth at \(j\). Conversely these data satisfy every
conjunct of \(\iota_+\). This proves the \(U_1\) clause.

The argument applies to the entire truth set of any operand
formula \(\alpha\), since the unary operations act on
that set. Integers \(j-1,j+1\) always exist in the fixed
order, even when the operand set has a first or last member.
If no strict occurrence exists in the relevant direction,
the outer existential is false. If one exists, its positive
integer distance from \(i\) has a least value. The
intervening falsehood condition says exactly that this nearest
occurrence is chosen, and the remaining clause checks its
neighbor on the far side. At distance one the intervening
interval is empty; at distance zero no witness is permitted.
No clause requires a value for the operand at the current
point itself. This verifies all stated boundary conditions
\cite[ch.~IV, def.~2, p.~100 and pp.~101--102]{Kamp1968}.
\end{proof}
\end{annexitem}

\begin{annexitem}{annex:p9.kamp.iv2.marker.lemma}{thm:p09-kamp-iv2-marker-lemma}{The complete successful, failed and absent endpoint cases}
\begin{proof}
Fix the arbitrary current point \(t\) and direction
\(\delta\in\{+1,-1\}\). For every integer \(m\),
write \(a_m\) and \(b_m\) for the Boolean truth values
of \(q_0,q_1\) at \(t+\delta m\), respectively.
Write \(c_m=\neg a_m\land b_m\). A unit shift
\(N_\delta\) increases \(m\) by one, while
\(N_{-\delta}\) decreases it by one. For \(\delta=+1\)
these are \(F_1,P_1\); for \(\delta=-1\) they are
\(P_1,F_1\). These identities follow from
\cref{thm:p09-kamp-iv2-unary-semantics} at every integer.

Let \(e_m\) be the truth of the directional marker
\(\eta_\delta\) at \(t+\delta m\). Expanding the
actual formulas in both rows of the directional table gives
\[
\begin{aligned}
 r_m&=a_m\land\neg a_{m-1}\land b_{m-1},\\
 k_m&=\neg a_m\land\neg b_m\land
                         \neg a_{m-1}\land\neg a_{m-2},\\
 e_m&=r_m\lor k_m\lor k_{m+1}.
\end{aligned}
 \tag{4}
\]
For example, the last disjunct is \(F_1B_P\) in
\(\eta_1\), and \(P_1B_F\) in \(\eta_2\);
both evaluate the failure condition one step farther along
the chosen search. The first two inspections in each
failure condition go toward the starting point. Thus (4)
retains every disjunct of the two source marker expressions.

The isolation clause in either direction becomes exactly
\[
 (J_\delta\eta_\delta)(t)
 \quad\Longleftrightarrow\quad
 \exists m\geq1\bigl(e_m\land\neg e_{m+1}\land
             \forall r(1\leq r<m\to\neg e_r)\bigr).
 \tag{5}
\]
For future isolation its far neighbor is \(j+1\); for
past isolation it is \(j-1\). Both are represented by
index \(m+1\) in (5). The interval between the witness
and the starting point is represented by the positive
indices less than \(m\).

By hypothesis \(c_1,c_2\) are true. In particular
\(a_1=a_2=\text{false}\) and
\(b_1=b_2=\text{true}\), so (4) gives \(e_1\)
false. We now exhaust all possible first exits from this
initial run of points satisfying \(C\).

\paragraph{A successful first exit.}
Suppose the least positive integer \(m\) for which
\(c_m\) fails has \(a_m\) true. Necessarily
\(m\geq3\). All positive earlier indices satisfy
\(\neg a_r\land b_r\), so \(m\) is a valid
strict witness for the target: its intervening \(b_r\)
are true. Its own \(b_m\) is unrestricted.

Here \(r_m\) is true, so \(e_m\) is true. For
\(1\leq r<m\), the first two disjuncts in (4) are
false because \(a_r\) is false and \(b_r\) true.
The last disjunct \(k_{r+1}\) is false because
\(b_{r+1}\) is true when \(r+1<m\), and because
\(a_m\) is true when \(r+1=m\). Hence no earlier
positive marker occurs. At \(m+1\), the term
\(r_{m+1}\) is false because it requires
\(\neg a_m\). The term \(k_{m+1}\) has the same
false requirement, and \(k_{m+2}\) also requires
\(\neg a_m\). Thus \(e_{m+1}\) is false, regardless
of every predicate value beyond the successful endpoint.
Equation (5) is therefore true.

\paragraph{A failed first exit.}
Suppose instead the least positive exit \(m\) has
\(a_m\) false. Since \(c_m\) fails, \(b_m\) must
also be false, and \(m\geq3\). There is no earlier
target witness because all earlier \(a_r\) are false.
The point at index \(m\) is not a witness either, and
its false \(b_m\) prevents every later index from
being a target witness. Thus the target existential fails.

The two preceding \(a\)-values are false, so
\(k_m\) is true. Equation (4) consequently makes
both \(e_{m-1}\) and \(e_m\) true. For
\(1\leq r<m-1\), each of \(r_r,k_r,k_{r+1}\)
is false: the positive indices \(r,r+1\) are still
in the run where \(a\) is false and \(b\) true.
Therefore the nearest positive marker is at \(m-1\),
and its far neighbor \(m\) is also a marker. It fails
the isolation requirement in (5). No later marker can
serve as the witness in (5), since the marker at
\(m-1\) would then be intervening. Hence (5) is false.

\paragraph{No exit.}
If no positive index fails \(C\), then every positive
\(a_r\) is false and every positive \(b_r\) is true.
All three disjuncts of \(e_r\) are false for every
positive \(r\), by (4). There is no marker witness
in (5), and no target witness either.

These cases exhaust all interpretations: any nonempty set
of positive exit indices has a least member. A successful
exit gives both existentials; a failed exit or no exit gives
neither. Equivalently, if a target witness exists at some
index, the first exit occurs no later than that witness,
and cannot be a failed exit because its false interval
value would lie before the witness. This proves both
directions of the claimed equivalence.

\paragraph{The source's no-witness alternative.}
One can also verify directly the case where every positive
\(a_r\) is false, allowing any later values of \(b_r\).
This is the source's separate alternative on pp.~104--105.
All positive \(r_r\) vanish, so
\(e_r=k_r\lor k_{r+1}\). We already know \(e_1\)
is false. If no positive marker exists, (5) is false.
Otherwise let \(h\geq2\) be its least positive index.
The term \(k_h\) cannot be true: it would make
\(e_{h-1}\) true, with \(h-1\geq1\), contradicting
minimality. Hence \(k_{h+1}\) is true, making
\(e_{h+1}\) true. The nearest marker again has a
true far neighbor and cannot satisfy (5). This supplies
the entire no-witness branch, including an unbounded tail
of failures or no failures at all.

All the calculations were made at arbitrary \(t\), using
both values of \(\delta\). In the past direction a
successful endpoint \(j=t-m\) is the nearest marker
and has \(\eta_2(j-1)\) false. A failed endpoint
\(j=t-m\) gives the adjacent markers at \(j+1,j\),
and the nearest one's farther, earlier neighbor is true.
No past exit means no past marker. In the future direction
the corresponding points are \(j=t+m\), its neighbor
\(j+1\), and the failed pair \(j-1,j\). This
establishes the two orientations with their actual operand
and neighbor order, including every case of the source
construction
\cite[ch.~IV, thm.~2 proof, pp.~101--105]{Kamp1968}.
\end{proof}
\end{annexitem}

\begin{annexitem}{annex:p9.kamp.iv2}{thm:p09-kamp-iv2}{Both explicit definitions at every integer position}
\begin{proof}
Fix \(t\in\mathbb Z\), a valuation of the two input
letters, and a direction \(\delta\in\{+1,-1\}\).
Use the Boolean values \(a_m,b_m,c_m\) and directional
marker of the preceding proof. Expanding either of the
two displayed target formulas gives a conjunction
\(g\land \ell\), where
\[
\begin{aligned}
 g&=a_1\lor\bigl(b_1\land(a_2\lor b_2)\bigr),\\
 \ell&=a_1\lor(b_1\land a_2)\lor
              \bigl(c_1\land c_2\land
                    (J_\delta\eta_\delta)(t)\bigr).
\end{aligned}
 \tag{6}
\]
For \(\delta=+1\) this is precisely \(\zeta_1\),
and for \(\delta=-1\) it is precisely \(\zeta_2\).
The interpretation uses two shifts in the same direction
where index two occurs. The desired target assertion in
the chosen direction is
\[
 \exists m\geq1\bigl(a_m\land
                  \forall r(1\leq r<m\to b_r)\bigr).
 \tag{7}
\]
We compare (6) and (7) in all possible initial cases.

If \(a_1\) is true, both \(g\) and \(\ell\) are
true by their first disjuncts. Index one witnesses (7),
with no intervening index and no condition on \(b_1\).
Suppose henceforth that \(a_1\) is false. If \(b_1\)
is false, index one is not a witness and every later
witness would have that failed interval point before it.
Thus (7) is false. All disjuncts of \(\ell\) are
false as well, because \(c_1\) requires \(b_1\).

Now assume \(a_1\) false and \(b_1\) true. If
\(a_2\) is true, the first conjunct \(g\) and the
second disjunct of \(\ell\) are true. Index two is
a witness for (7), whose only intervening value is
\(b_1\). This case imposes no requirement on \(b_2\).
If instead \(a_2\) and \(b_2\) are both false,
neither initial index is a witness and the failed
\(b_2\) prevents all later witnesses. Both (7) and
\(\ell\), whose third disjunct requires \(c_2\),
are false. The guard \(g\) is false too.

The only remaining case is
\(\neg a_1\land b_1\land\neg a_2\land b_2\),
that is, \(c_1\land c_2\). Here \(g\) is true,
and \(\ell\) reduces exactly to
\((J_\delta\eta_\delta)(t)\). Neither index one
nor index two can be a witness for (7), so (7) reduces
to its restriction \(m\geq3\). The equivalence of
these two remaining conditions is exactly
\cref{thm:p09-kamp-iv2-marker-lemma}. This completes
the exhaustive comparison in both directions.

Finally, for \(\delta=+1\), the bijection
\(m\mapsto s=t+m\) identifies positive indices with
all integers \(s>t\). The indices \(1\leq r<m\)
are exactly the integers in \((t,s)\). Thus (7) is
the source's strict \(U(A,B)\) assertion. For
\(\delta=-1\), the bijection \(m\mapsto s=t-m\)
identifies positive indices with \(s<t\), and the
points \(t-r\), \(1\leq r<m\), are exactly the
integers in \((s,t)\). Hence (7) is strict
\(S(A,B)\). The input \(a_m\) always belongs to
the witness operand, and \(b_r\) to the open-interval
operand. Since \(t\) and the valuation were arbitrary,
these are the two operation identities at every integer.

Every abbreviation in the construction expands into
\(S_1,U_1,P_1,F_1,\neg,\land\). Disjunction is
\(\neg(\neg\alpha\land\neg\beta)\), and the
auxiliary formulas are finite expressions using only
\(q_0,q_1\). Consequently the displayed formulas
express the targets over exactly the source's six-operation
basis, without importing ordinary unbounded past or future.

\paragraph{Boundary examples.}
For \(A=\{3,4,5,\ldots\},B=\mathbb Z,t=0\),
the failure condition \(B_P\) is everywhere false,
and \(q_0\land P_1\neg q_0\land P_1q_1\) is
true exactly at three. Thus \(\eta_1\) has that
single positive occurrence, with its far neighbor false.
Both initial points satisfy \(C\), so the hard
disjunct of \(\zeta_1\) is true. There is no need
for the witness proposition to stop being true after
the endpoint. Replacing every integer by its negative
interchanges \(P_1,F_1\), \(S_1,U_1\), and the
two marker and target formulas by their proved clauses.
This gives the stated past example, with its endpoint
at minus three and its farther neighbor at minus four.

For \(A=\varnothing,B=\{1,2\}\), the first failed
future point is three. The term \(B_P\) is true
there, so \(F_1B_P\) is true at two. Hence both
two and three are markers and the future isolation
condition fails. For \(A=\varnothing,B=\mathbb Z\),
both the success and failure conditions vanish at all
positive points, so there is no future marker. These
are the failed and absent cases of the marker lemma.

For \(A=\{1\}\), the first initial case in (6)
applies for every \(B\), including the empty set.
For \(A=\{2\}\), the second-witness case applies
exactly when \(1\in B\); if that value is false
it blocks every possible witness. Reflecting gives
the identical statements for witnesses minus one and
minus two in the past. A value of the witness
proposition at the current point alone cannot supply
a strict witness. In particular \(A=\{0\}\),
\(B=\{-1,-2\}\) makes \(\zeta_2(0)\) false:
there is no past witness and the no-witness marker
calculation applies.

If the witness set is empty, (7) is false in both
directions; if it is all integers, its immediate
neighbor supplies a witness. If the interval set is
empty, only an immediate witness can work; if it is
all integers, (7) asks for any strict witness in the
chosen direction. The preceding equivalence includes
all these cases and every bounded or unbounded input
set. The order itself is exactly \(\mathbb Z\),
with no first or last time point; endpoints of predicate
blocks and vacuous open intervals have all been handled
\cite[ch.~IV, thm.~2, pp.~100--105]{Kamp1968}.
\end{proof}
\end{annexitem}
